\documentclass[10pt,a4paper]{amsart}
\usepackage[margin=19mm]{geometry}
\usepackage{amsmath,amssymb,mathtools}
\usepackage[scr=rsfs,cal=euler]{mathalfa}
\usepackage{xfrac}
\usepackage[unicode,hidelinks,bookmarks=false]{hyperref}
\newcommand{\ie}{i.e.\ }
\newcommand{\cf}{cf.\ }
\newcommand{\resp}{respectively\ }
\newcommand{\Subsection}{\subsection}
\providecommand{\nobreakdash}{\nobreak\hbox{-}}
\setlength{\emergencystretch}{2em}
\multlinegap0pt
\usepackage{amssymb}
\usepackage[shortlabels]{enumitem}
\usepackage{dsfont}
\newtheorem{lemma}{Lemma}
\newcommand{\dn}{\mathord{\downarrow}\hspace{0.05em}}
\title{A new method for constructing a compact and non SI-compact space}
\author{Zhengmao He, Bin Zhao}
\date{Source: arXiv:2608.01048v1 (CC BY 4.0)}
\begin{document}
\maketitle

\noindent\emph{Source and licence.} A new method for constructing a compact and non SI-compact space. Version arXiv:2608.01048v1 (CC BY 4.0), distributed by arXiv under the Creative Commons Attribution 4.0 International licence, http://creativecommons.org/licenses/by/4.0/, as recorded in the arXiv licence metadata for this exact version and shown on the abstract page https://arxiv.org/abs/2608.01048v1. Full licence notice: \texttt{licenses/arxiv-2608.01048-CC-BY-4.0.txt}.\par\smallskip

\noindent\textit{Editorial scope.} Counted unit: the article's lemma lemma (source0.tex lines 746--754). The article's complete original proof of it is delivered with it (source0.tex lines 756--945, one \texttt{proof} environment of 190 lines). No mathematics is written, repaired or re-derived: the statement and the proof are the article's own lines, in the article's own notation, and no retained line is substituted or re-flowed.  No further source-local context is needed: every cross-reference of every retained line resolves inside the two delivered spans.  Nothing was omitted from any retained span, so the package carries no omission record.  No retained line contains a citation, so the wrapper carries no bibliography and no bibliography entry.  No retained line contains a figure environment, an image-inclusion command or drawing code, so no image is delivered and the package declares no asset dependency.  Editorial formatting only, in addition: an AMS \texttt{amsart} preamble that restates the article's own declarations and macros used by the retained lines, namely the environments \texttt{lemma} on separate counters, together with the standard packages those lines need.  The retained environments print in this document as Lemma 1 in order of appearance, whereas the article prints the same items with the numbering of its own full document.  The complete native source of the article as distributed in the arXiv e-print for this exact version is bundled unchanged as \texttt{sources/206590/source0.tex}.
\begin{lemma} {\rm Let $A\subseteq\widehat{P}$ be subset such that

(1) every finite subset of $A$ has an upper bound.

(2) $A$ just meet $L_{n_{0}}$, $L_{n_{1}}$ and  $L_{n_{2}}$ nonempty.

Then $A$ has an upper bound.}

\end{lemma}

\begin{proof}We need only to consider the case that $A$ is infinite. Assume that $A\cap L_{n_{0}}$ is infinite and $A\cap L_{n_{1}}\neq\emptyset$, $A\cap L_{n_{2}}\neq\emptyset$. There are only 8 possibilities ((1)-(8)).

(1) $A\cap(\{n_{0}\}\times\mathbb{N})$ is infinite and $A\cap(\{n_{1}\}\times\mathbb{N})\neq\emptyset$, $A\cap(\{n_{2}\}\times\mathbb{N})\neq\emptyset$.

Choose a $(n_{i},m_{i})\in A\cap(\{n_{i}\}\times\mathbb{N})$, $i=1,2,3$.  By strict orders $\prec_{2},\prec_{3} $ and $\prec_{4}$,\\ $\{(n_{1},m_{1}),(n_{2},m_{2}),(n_{3},m_{3})\}$ has no upper bounds. So case (1) is impossible.

(2) $A\cap(\{n_{0}\}\times\mathbb{N})$ is infinite and $A\cap(\{n_{1}\}\times\mathbb{N}^{<\omega})\neq\emptyset$, $A\cap(\{n_{2}\}\times\mathbb{N})\neq\emptyset$.

Choose a $(n_{0},m_{0})\in A\cap(\{n_{0}\}\times\mathbb{N})$, a $(n_{1},\xi_{1}\xi_{2}\cdot\cdot\cdot\xi_{m})\in A\cap(\{n_{1}\}\times\mathbb{N}^{<\omega})$ and a $(n_{2},m_{2})\in A\cap(\{n_{2}\}\times\mathbb{N})$.

Then $\{(n_{0},m_{0}),(n_{1},\xi_{1}\xi_{2}\cdot\cdot\cdot\xi_{m}),(n_{2},m_{2})\}$ has an upper bound $l$. As $n_{0}\neq n_{2}$, $l=(n_{0},\top)$ or $l=(n_{2},\top)$.



(2.1) $l=(n_{2},\top)$.

Since $A\cap(\{n_{0}\}\times\mathbb{N})$ is infinite, the element $m_{0}$ has infinitely many possibilities. Hence $l=(n_{2},\top)$ is impossible.



%we can take a $(n_{0},m'_{0})\in A$ such that $$m'_{0}>max\{m_{0}+2,\xi_{1}+2,\xi_{2}+2,\cdot\cdot\cdot,\xi_{m}+2,m_{2}+2\}.$$
%Clearly $(n_{1},\xi_{1}\xi_{2}\cdot\cdot\cdot\xi_{m})\prec_{2}j=(n_{2},\top)$ and $(n_{0},m_{0})\leq j=(n_{2},\top)$. By $(n_{1},\xi_{1}\xi_{2}\cdot\cdot\cdot\xi_{m})\prec_{2}j=(n_{2},\top)$, we have $n_{2}=f_{n_{1}r_{1}}(\xi'_{1}\xi'_{2}\cdot\cdot\cdot\xi'_{m'})$ for some $n_{1}<r_{1}$ and $\xi_{1}\xi_{2}\cdot\cdot\cdot\xi_{m}\leq\xi'_{1}\xi'_{2}\cdot\cdot\cdot\xi'_{m'}$ in $\mathbb{N}^{<\omega}$.

%(2.1.1) If $m'=1$, then $m=1$ and $(n_{0},m_{0})\prec_{3} j=(n_{2},\top)$ with $\xi_{1}\xi_{2}\cdot\cdot\cdot\xi_{m}=\xi'_{1}\xi'_{2}\cdot\cdot\cdot\xi'_{m'}=\xi_{1}$. Whence, $(n_{0},m_{0})\prec_{3}=(f_{n_{1}r_{1}}(\xi_{1}),\top)$ and then $n_{0}=r_{1}$ and $m_{0}\leq\xi_{1}$. Now, we have $n_{0}=r_{1}>n_{1}$, $m_{0}\leq\xi_{1}$ and $n_{2}=f_{n_{1}r_{1}}(\xi_{1})$. In addition, $\{(n_{0},m'_{0}),(n_{1},\xi_{1}),(n_{2},m_{2})\}$ also has an upper bound $j'$. Similarly, $j'=(n_{0},\top)$ or $j'=(n_{2},\top)$. we consider the following two possibilities ((2.1.1-1)
 %and (2.1.1-2)).

%(2.1.1-1) $j'=(n_{0},\top)$.

%Then $(n_{1},\xi_{1})\prec_{2}(n_{0},\top)$. This means that $n_{0}=f_{n_{1}r_{2}}(\xi''_{1}\xi''_{2}\cdot\cdot\cdot\xi''_{m''})$ for some $n_{1}<r_{2}$ and $\xi_{1}\leq\xi''_{1}\xi''_{2}\cdot\cdot\cdot\xi''_{m''}$ in $\mathbb{N}^{<\omega}$.

%If $m''=1$, then $(n_{2},m_{2})\prec_{3}j'=(n_{0},\top)$ and $\xi_{1}=\xi''_{1}\xi''_{2}\cdot\cdot\cdot\xi''_{m''}$ . Thus, $n_{2}=r_{2}$ and $m_{2}\leq\xi_{1}$. As a result, $n_{2}=f_{n_{1}r_{1}}(\xi_{1})=f_{n_{1}n_{0}}(\xi_{1})>n_{0}$ and $n_{0}=f_{n_{1}r_{2}}(\xi_{1})=f_{n_{1}n_{2}}(\xi_{1})>n_{2}$, impossible.

%If $1<m''$, then $(n_{2},m_{2})\prec_{4}j'=(n_{0},\top)$. Thus $n_{2}=f_{n_{1}r_{2}}(\xi''_{1}\xi''_{2}\cdot\cdot\cdot\xi''_{m''-1})$ and $m_{2}\leq \xi''_{m''}$. Now, we have $n_{0}=f_{n_{1}r_{2}}(\xi''_{1}\xi''_{2}\cdot\cdot\cdot\xi''_{m''})>f_{n_{1}r_{2}}(\xi''_{1}\xi''_{2}\cdot\cdot\cdot\xi''_{m''-1})=n_{2}$ and $n_{2}=f_{n_{1}r_{1}}(\xi_{1})=f_{n_{1}n_{0}}(\xi_{1})>n_{0}$, impossible.

%(2.1.1-2) $j'=(n_{2},\top)$.

%Then $(n_{1},\xi_{1})\prec_{2}(n_{2},\top)$. This leads to $n_{2}=f_{n_{1}r_{3}}(\overline{\xi}_{1}\overline{\xi}_{2}\cdot\cdot\cdot\overline{\xi}_{\overline{m}})$ for some $n_{1}<r_{3}$ and $\xi_{1}\leq\overline{\xi}_{1}\overline{\xi}_{2}\cdot\cdot\cdot\overline{\xi}_{\overline{m}}$ in $\mathbb{N}^{<\omega}$.


%If $\overline{m}=1$, then $(n_{0},m'_{0})\prec_{3}j'=(n_{2},\top)$ and $\xi_{1}=\overline{\xi}_{1}\overline{\xi}_{2}\cdot\cdot\cdot\overline{\xi}_{\overline{m}}$ . Thus, $n_{2}=r_{3}$ and $m'_{0}\leq\xi_{1}$, impossible.

%If $1<m''$, then $(n_{0},m'_{0})\prec_{4}j'=(n_{2},\top)$. Thus $n_{0}=f_{n_{1}r_{3}}(\overline{\xi}_{1}\overline{\xi}_{2}\cdot\cdot\cdot\overline{\xi}_{\overline{m}-1})$ and $m'_{0}\leq \overline{\xi}_{\overline{m}}$. Now, we have $n_{2}=f_{n_{1}r_{3}}(\overline{\xi}_{1}\overline{\xi}_{2}\cdot\cdot\cdot\overline{\xi}_{\overline{m}})=f_{n_{1}r_{1}}(\xi_{1})=f_{n_{1}n_{0}}(\xi_{1})$.
%Thus, we have that $n_{0}=r_{1}=r_{3}=f_{n_{1}r_{3}}(\overline{\xi}_{1}\overline{\xi}_{2}\cdot\cdot\cdot\overline{\xi}_{\overline{m}-1})>r_{3}$, a contradiction.





%(2.1.2) If $1<m'$, then $(n_{0},m_{0})\prec_{4} j=(n_{2},\top)$. Hence, $n_{0}=f_{n_{1}r_{1}}(\xi'_{1}\xi'_{2}\cdot\cdot\cdot\xi'_{m'-1})$ and $m_{0}\leq\xi'_{m'}$. Now, we have $n_{0}=f_{n_{1}r_{1}}(\xi'_{1}\xi'_{2}\cdot\cdot\cdot\xi'_{m'-1})$ and $m_{0}\leq\xi'_{m'}$, $n_{2}=f_{n_{1}r_{1}}(\xi'_{1}\xi'_{2}\cdot\cdot\cdot\xi'_{m'})$. Note that $\{(n_{0},m''_{0}),(n_{1},\xi_{1}\xi_{2}\cdot\cdot\cdot\xi_{m}),(n_{2},m_{2})\}$ also has an upper bound $j_{1}'$, where $(n_{0},m''_{0})\in A$ with $m''_{0}>max\{m'_{0}+1,\xi'_{1}+1,\xi'_{2}+1,\cdot\cdot\cdot,\xi'_{m'}+1\}$. Similarly, $j_{1}'=(n_{0},\top)$ or $j_{1}'=(n_{2},\top)$. we just consider the following two possibilities ((2.1.2-1) and (2.1.2-2)).

%(2.1.2-1) $j_{1}'=(n_{0},\top)$.

% Then $(n_{1},\xi_{1}\xi_{2}\cdot\cdot\cdot\xi_{m})\prec_{2}(n_{0},\top)$. This means that $n_{0}=f_{n_{1}r_{4}}(\xi'''_{1}\xi'''_{2}\cdot\cdot\cdot\xi'''_{m'''})$ for some $n_{1}<r_{4}$ and $\xi_{1}\xi_{2}\cdot\cdot\cdot\xi_{m}\leq\xi'''_{1}\xi'''_{2}\cdot\cdot\cdot\xi'''_{m'''}$ in $\mathbb{N}^{<\omega}$.

%If $m'''=1$, then $(n_{2},m_{2})\prec_{3}j'=(n_{0},\top)$ and $\xi_{1}=\xi'''_{1}\xi'''_{2}\cdot\cdot\cdot\xi'''_{m'''}$. Thus, $n_{2}=r_{4}$ and $m_{2}\leq\xi_{1}$. As a result, $n_{2}=f_{n_{1}r_{1}}(\xi'_{1}\xi'_{2}\cdot\cdot\cdot\xi'_{m'})>f_{n_{1}r_{1}}(\xi'_{1}\xi'_{2}\cdot\cdot\cdot\xi'_{m'-1})=n_{0}$ and $n_{0}=f_{n_{1}r_{4}}(\xi_{1})=f_{n_{1}n_{2}}(\xi_{1})>n_{2}$, impossible.

%If $1<m''$, then $(n_{2},m_{2})\prec_{4}j'=(n_{0},\top)$. Thus $n_{2}=f_{n_{1}r_{4}}(\xi'''_{1}\xi'''_{2}\cdot\cdot\cdot\xi'''_{m'''-1})$ and $m_{2}\leq \xi'''_{m'''}$. Consequently, we have that $n_{0}=f_{n_{1}r_{4}}(\xi'''_{1}\xi'''_{2}\cdot\cdot\cdot\xi'''_{m'''})>f_{n_{1}r_{4}}(\xi'''_{1}\xi'''_{2}\cdot\cdot\cdot\xi'''_{m'''-1})=n_{2}$ and  $n_{2}=f_{n_{1}r_{1}}(\xi'_{1}\xi'_{2}\cdot\cdot\cdot\xi'_{m'})>f_{n_{1}r_{1}}(\xi'_{1}\xi'_{2}\cdot\cdot\cdot\xi'_{m'-1})=n_{0}$, impossible.

%(2.1.2-2) $j_{1}'=(n_{2},\top)$.

%Then $(n_{1},\xi_{1}\xi_{1}\xi_{2}\cdot\cdot\cdot\xi_{m})\prec_{2}(n_{2},\top)$. This leads to $n_{2}=f_{n_{1}r_{5}}(\widetilde{\xi}_{1}\widetilde{\xi}_{2}\cdot\cdot\cdot\widetilde{\xi}_{\widetilde{m}})$ for some $n_{1}<r_{5}$ and $\xi_{1}\xi_{2}\cdot\cdot\cdot\xi_{m}\leq\widetilde{\xi}_{1}\widetilde{\xi}_{2}\cdot\cdot\cdot\widetilde{\xi}_{\widetilde{m}}$ in $\mathbb{N}^{<\omega}$.


%If $\widetilde{m}=1$, then $(n_{0},m''_{0})\prec_{3}j_{1}'=(n_{2},\top)$ and $\xi_{1}=\widetilde{\xi}_{1}\widetilde{\xi}_{2}\cdot\cdot\cdot\widetilde{\xi}_{\widetilde{m}}$. Hence, $n_{2}=r_{5}$ and $m''_{0}\leq\xi_{1}$, impossible.

%If $1<\widetilde{m}$, then $(n_{0},m'_{0})\prec_{4}j_{1}'=(n_{2},\top)$. Thus $n_{0}=f_{n_{1}r_{5}}(\widetilde{\xi}_{1}\widetilde{\xi}_{2}\cdot\cdot\cdot\widetilde{\xi}_{\widetilde{m}-1})$ and $m''_{0}\leq \widetilde{\xi}_{\widetilde{m}}$. Now, we have $n_{2}=f_{n_{1}r_{5}}(\widetilde{\xi}_{1}\widetilde{\xi}_{2}\cdot\cdot\cdot\widetilde{\xi}_{\widetilde{m}})=f_{n_{1}r_{1}}(\xi'_{1}\xi'_{2}\cdot\cdot\cdot\xi'_{m'})$. Hence, we have $\widetilde{\xi}_{1}\widetilde{\xi}_{2}\cdot\cdot\cdot\widetilde{\xi}_{\widetilde{m}}=\xi'_{1}\xi'_{2}\cdot\cdot\cdot\xi'_{m'}$ and then $\widetilde{\xi}_{\widetilde{m}}=\xi'_{m'}\geq m''_{0}$, impossible.

(2.2) $l=(n_{0},\top)$

Then $(n_{1},\xi_{1}\xi_{2}\cdot\cdot\cdot\xi_{m})\prec_{2}l=(n_{0},\top)$ and $(n_{2},m_{2})\leq l=(n_{0},\top)$. By $(n_{1},\xi_{1}\xi_{2}\cdot\cdot\cdot\xi_{m})\prec_{2}l=(n_{0},\top)$, $n_{0}=f_{n_{1}r_{1}}(\alpha_{1}\alpha_{2}\cdot\cdot\cdot\alpha_{t})$ for some $n_{1}<r_{1}$ and $\xi_{1}\xi_{2}\cdot\cdot\cdot\xi_{m}\leq\alpha_{1}\alpha_{2}\cdot\cdot\cdot\alpha_{t}$ in $\mathbb{N}^{<\omega}$. There are only two cases (2.2.1) and (2.2.2).

(2.2.1) $t=1$.

In this case, $\xi_{1}\xi_{2}\cdot\cdot\cdot\xi_{m}=\alpha_{1}\alpha_{2}\cdot\cdot\cdot\alpha_{t}=\xi_{1}$ and $(n_{2},m_{2})\prec_{3}\leq l=(n_{0},\top)=(f_{n_{1}r_{1}}(\xi_{1}),\top)$. Then $n_{2}=r_{1}>n_{1}$ and $m_{2}\leq\xi_{1}$. In the following, we assert that $$A\subseteq\dn(n_{0},\top)=\dn(f_{n_{1}r_{1}}(\xi_{1}),\top).$$
Assume that $A\not\subseteq\dn(f_{n_{1}r_{1}}(\xi_{1}),\top)$. Take a $(n_{3},m_{3})\in A\setminus\dn(f_{n_{1}r_{1}}(\xi_{1}),\top)$. Note that
$$\dn(f_{n_{1}r_{1}}(\xi_{1}),\top)=L_{f_{n_{1}r_{6}}(\xi_{1})}\cup\{(n_{1},\xi_{1})\}\cup\{(r_{1},k)\mid 1\leq k\leq\xi_{1}\},$$
where $(n_{1},\xi_{1})\in\{n_{1}\}\times\mathbb{N}^{<\omega}$. So we just consider the following 4 cases (2.2.1-1)-(2.2.1-4).

(2.2.1-1) $n_{3}=n_{1}$, $m_{3}\in\mathbb{N}$.

Then $\{(n_{0},m_{0}),(n_{1},\xi_{1}),(n_{2},m_{2}),(n_{1},m_{3})\}$ has an upper bound $j_{2}$. This is impossible since $n_{0},n_{1},n_{2}$ are different.

(2.2.1-2) $n_{3}=n_{1}$, $m_{3}\in\mathbb{N}^{<\omega}\setminus\{\xi_{1}\}$.

Let $m_{3}=\beta_{1}\beta_{2}\cdot\cdot\cdot\beta_{t_{1}}$. Then $\{(n_{0},m_{0}),(n_{1},\xi_{1}),(n_{2},m_{2}),(n_{1},\beta_{1}\beta_{2}\cdot\cdot\cdot\beta_{t_{1}})\}$ has an upper bound $j'_{2}$.
Similarly, $j'_{2}=(n_{0},\top)$ or $j'_{2}=(n_{2},\top)$.

If $j'_{2}=(n_{0},\top)$, $(n_{1},\beta_{1}\beta_{2}\cdot\cdot\cdot\beta_{t_{1}})\prec_{2}(n_{0},\top)$. This implies $n_{0}=f_{n_{1}r_{2}}(\beta'_{1}\beta'_{2}\cdot\cdot\cdot\beta'_{t'_{1}})$ for some $n_{1}<r_{2}$ and $\beta_{1}\beta_{2}\cdot\cdot\cdot\beta_{t_{1}}\leq\beta'_{1}\beta'_{2}\cdot\cdot\cdot\beta'_{t'_{1}}$. Now, we have $n_{0}=f_{n_{1}r_{2}}(\beta'_{1}\beta'_{2}\cdot\cdot\cdot\beta'_{t'_{1}})=f_{n_{1}r_{1}}(\xi_{1})$ and then $\beta'_{1}\beta'_{2}\cdot\cdot\cdot\beta'_{t'_{1}}=\xi_{1}$. Since $\beta_{1}\beta_{2}\cdot\cdot\cdot\beta_{t_{1}}\leq\beta'_{1}\beta'_{2}\cdot\cdot\cdot\beta'_{t'_{1}}=\xi_{1}$, we have $\beta_{1}\beta_{2}\cdot\cdot\cdot\beta_{t_{1}}=\xi_{1}$, a contradiction.

If $j'_{2}=(n_{2},\top)$, then $(n_{0},m_{0})\prec_{3}(n_{2},\top)$ or $(n_{0},m_{0})\prec_{4}(n_{2},\top)$. Hence, by the strict orders $\prec_{3}$ and $\prec_{4}$, $n_{0}<n_{2}$. However, $n_{2}=r_{1}<f_{n_{1}r_{1}}(\xi_{1})=n_{0}$, a contradiction.

(2.2.1-3) $n_{3}=n_{2}=r_{1}$, $m_{3}\in\mathbb{N}^{<\omega}$.

Let $m_{3}=\beta'_{1}\beta'_{2}\cdot\cdot\cdot\beta'_{t'_{1}}$. Then $\{(n_{0},m_{0}),(n_{1},\xi_{1}),(n_{2},m_{2}),(n_{2},\beta'_{1}\beta'_{2}\cdot\cdot\cdot\beta'_{t'_{1}})\}$ has an upper bound $j_{3}$. As $n_{1}\neq n_{2}$ and $\xi_{1},\beta'_{1}\beta'_{2}\cdot\cdot\cdot\beta'_{t'_{1}}\in\mathbb{N}^{<\omega}$, $j_{3}=(n_{1},\top)$ or $j_{3}=(n_{2},\top)$.

If $j_{3}=(n_{1},\top)$, $(n_{0},m_{0})\prec_{3}(n_{1},\top)$ or $(n_{0},m_{0})\prec_{4}(n_{1},\top)$. Using the strict orders $\prec_{3}$ and $\prec_{4}$, $n_{0}<n_{1}$. However, $n_{1}<f_{n_{1}r_{6}}(\xi_{1})=n_{0}$, a contradiction.

If $j_{3}=(n_{2},\top)$, $(n_{0},m_{0})\prec_{3}(n_{2},\top)$ or $(n_{0},m_{0})\prec_{4}(n_{2},\top)$. Similarly, $n_{0}<n_{2}$ but $n_{2}=r_{1}<f_{n_{1}r_{6}}(\xi_{1})=n_{0}$, a contradiction.

(2.2.1-4) $n_{3}=n_{2}=r_{1}$, $m_{3}\in\mathbb{N}$ with $\xi_{1}<m_{3}$.

Then $\{(n_{0},m_{0}),(n_{1},\xi_{1}),(n_{2},m_{3})\}$ has an upper bound $j'_{3}$. Similarly, $j'_{3}=(n_{0},\top)$ or $j'_{3}=(n_{2},\top)$.

If $j'_{3}=(n_{2},\top)$, $(n_{0},m_{0})\prec_{3}(n_{2},\top)$ or $(n_{0},m_{0})\prec_{4}(n_{2},\top)$. So we have $n_{0}<n_{2}$ but $n_{2}=r_{1}<f_{n_{1}r_{1}}(\xi_{1})=n_{0}$, a contradiction.

If $j'_{3}=(n_{0},\top)$, then  $(n_{2},m_{3})\prec_{3}(n_{0},\top)=(f_{n_{1}r_{1}}(\xi_{1}),\top)$. Thus, $m_{3}\leq\xi_{1}$, a contradiction.


(2.2.2) $1<t$.

In this case $(n_{2},m_{2})\prec_{4}l=(n_{0},\top)=(f_{n_{1}r_{1}}(\alpha_{1}\alpha_{2}\cdot\cdot\cdot\alpha_{t}),\top)$. So we have that $n_{2}=f_{n_{1}r_{1}}(\alpha_{1}\alpha_{2}\cdot\cdot\cdot\alpha_{t-1})$ and $m_{2}\leq\alpha_{t}$. In the following, we assert that $$A\subseteq\dn(n_{0},\top)=\dn(f_{n_{1}r_{1}}(\alpha_{1}\alpha_{2}\cdot\cdot\cdot\alpha_{t}),\top).$$
Assume that $A\not\subseteq\dn(f_{n_{1}r_{1}}(\alpha_{1}\alpha_{2}\cdot\cdot\cdot\alpha_{t}),\top)$. Choose a $(n_{4},m_{4})\in A\setminus\dn(f_{n_{1}r_{1}}(\alpha_{1}\alpha_{2}\cdot\cdot\cdot\alpha_{t}),\top)$. Note that
\begin{center}$\begin{array}{lll}
\dn(f_{n_{1}r_{1}}(\alpha_{1}\alpha_{2}\cdot\cdot\cdot\alpha_{t}),\top)&=&L_{f_{n_{1}r_{1}}(\alpha_{1}\alpha_{2}\cdot\cdot\cdot\alpha_{t})}\cup\{(n_{1},\alpha_{1}\alpha_{2}\cdot\cdot\cdot\alpha_{t'})\mid 1\leq t'\leq t\} \\
&&\cup\{(f_{n_{1}r_{1}}(\alpha_{1}\alpha_{2}\cdot\cdot\cdot\alpha_{t-1}),k)\mid 1\leq k\leq\alpha_{t}\}.\\
\end{array}$\end{center}
So we just consider the following 4 cases (2.2.2-1)-(2.2.2-4).

(2.2.2-1) $n_{4}=n_{1}$, $m_{4}\in\mathbb{N}$.

Then
$$\{(n_{0},m_{0}),(n_{1},\xi_{1}\xi_{2}\cdot\cdot\cdot\xi_{m}),(n_{2},m_{2}),(n_{1},m_{4})\}$$ has an upper bound $j_{4}$. It is impossible since $n_{0},n_{1},n_{2}$ are different.

(2.2.2-2) $n_{4}=n_{1}$, $m_{4}\in\mathbb{N}^{<\omega}\setminus\{(n_{1},\alpha_{1}\alpha_{2}\cdot\cdot\cdot\alpha_{t'})\mid 1\leq t'\leq t\}$.

Let $m_{4}=\varphi_{1}\varphi_{2}\cdot\cdot\cdot\varphi_{t_{2}}$. Then $$\{(n_{0},m_{0}),(n_{1},\xi_{1}\xi_{2}\cdot\cdot\cdot\xi_{m}),(n_{2},m_{2}),(n_{1},\varphi_{1}\varphi_{2}\cdot\cdot\cdot\varphi_{t_{2}})\}$$ has an upper bound $j'_{4}$. Clearly, $j'_{4}=(n_{0},\top)$ or $j'_{4}=(n_{2},\top)$.

If $j'_{4}=(n_{2},\top)$, then $(n_{0},m_{0})\prec_{3}j'_{4}=(n_{2},\top)$ or $(n_{0},m_{0})\prec_{4}j'_{4}=(n_{2},\top)$. So we have $n_{0}<n_{2}$. However, $$n_{0}=f_{n_{1}r_{1}}(\alpha_{1}\alpha_{2}\cdot\cdot\cdot\alpha_{t})>f_{n_{1}r_{1}}(\alpha_{1}\alpha_{2}\cdot\cdot\cdot\alpha_{t-1})=n_{2},$$ a contradiction.

If $j'_{4}=(n_{0},\top)$, then $(n_{1},\xi_{1}\xi_{2}\cdot\cdot\cdot\xi_{m})\prec_{2}j'_{4}=(n_{0},\top)$ , $(n_{1},\varphi_{1}\varphi_{2}\cdot\cdot\cdot\varphi_{t_{2}})\prec_{2}j'_{4}=(n_{0},\top)$ and $(n_{2},m_{2})\leq(n_{0},\top)$.
Then $n_{0}=f_{n_{1}r_{3}}(\psi_{1}\psi_{2}\cdot\cdot\cdot\psi_{t_{3}})$ for some $n_{1}<r_{3}$ and $$\xi_{1}\xi_{2}\cdot\cdot\cdot\xi_{m}\leq\psi_{1}\psi_{2}\cdot\cdot\cdot\psi_{t_{3}}, \varphi_{1}\varphi_{2}\cdot\cdot\cdot\varphi_{t_{2}}\leq\psi_{1}\psi_{2}\cdot\cdot\cdot\psi_{t_{3}}. $$ Now, we have $n_{0}=f_{n_{1}r_{3}}(\psi_{1}\psi_{2}\cdot\cdot\cdot\psi_{t_{3}})=f_{n_{1}r_{6}}(\alpha_{1}\alpha_{2}\cdot\cdot\cdot\alpha_{t})$. Whence, $\psi_{1}\psi_{2}\cdot\cdot\cdot\psi_{t_{3}}=\alpha_{1}\alpha_{2}\cdot\cdot\cdot\alpha_{t}$ and then $\varphi_{1}\varphi_{2}\cdot\cdot\cdot\varphi_{t_{2}}\leq\alpha_{1}\alpha_{2}\cdot\cdot\cdot\alpha_{t}$, a contradiction.

(2.2.2-3) $n_{4}=n_{2}$, $m_{4}\in\mathbb{N}$ with $\alpha_{t}<m_{4}$.

Then $\{(n_{0},m_{0}),(n_{1},\xi_{1}\xi_{2}\cdot\cdot\cdot\xi_{m}),(n_{2},m_{4})\}$ has an upper bound $j_{5}$. Obviously, $j_{5}=(n_{0},\top)$ or $j_{5}=(n_{2},\top)$.

If $j_{5}=(n_{2},\top)$, then $(n_{0},m_{0})\prec_{3}j'_{4}=(n_{2},\top)$ or $(n_{0},m_{0})\prec_{4}j'_{4}=(n_{2},\top)$. Whence, $n_{0}<n_{2}$ but $$n_{0}=f_{n_{1}r_{1}}(\alpha_{1}\alpha_{2}\cdot\cdot\cdot\alpha_{t})>f_{n_{1}r_{1}}(\alpha_{1}\alpha_{2}\cdot\cdot\cdot\alpha_{t-1})=n_{2},$$ a contradiction.

If $j_{5}=(n_{0},\top)$, then $(n_{2},m_{4})\prec_{4}(n_{0},\top)=(f_{n_{1}r_{6}}(\alpha_{1}\alpha_{2}\cdot\cdot\cdot\alpha_{t}),\top)$. Then $m_{4}\leq\alpha_{t}$, impossible.

(2.2.2-4) $n_{4}=n_{2}$, $m_{4}\in\mathbb{N}^{<\omega}$.

Then $\{(n_{0},m_{0}),(n_{1},\xi_{1}\xi_{2}\cdot\cdot\cdot\xi_{m}),(n_{2},m_{2}),(n_{2},m_{4})\}$ has an upper bound $j'_{5}$. Since $n_{0}\neq n_{2}$ and $m_{0}, m_{2}\in\mathbb{N}$, $j'_{5}=(n_{0},\top)$ or $j'_{5}=(n_{2},\top)$. Simultaneously, as  $n_{1}\neq n_{2}$ and $\xi_{1}\xi_{2}\cdot\cdot\cdot\xi_{m}, m_{4}\in\mathbb{N}^{<\omega}$, $j'_{5}=(n_{1},\top)$ or $j'_{5}=(n_{2},\top)$. Consequently, $j'_{5}=(n_{2},\top)$ and $(n_{0},m_{0})\prec_{3}j'_{5}=(n_{2},\top)$ or $(n_{0},m_{0})\prec_{4}j'_{5}=(n_{2},\top)$. Similarly, we have $n_{0}<n_{2}$ but it contradicts  $n_{0}=f_{n_{1}r_{1}}(\alpha_{1}\alpha_{2}\cdot\cdot\cdot\alpha_{t})>f_{n_{1}r_{1}}(\alpha_{1}\alpha_{2}\cdot\cdot\cdot\alpha_{t-1})=n_{2}$, impossible.

So we can conclude that $A$ has an upper bound in case (2).



(3) $A\cap(\{n_{0}\}\times\mathbb{N})$ is infinite and $A\cap(\{n_{1}\}\times\mathbb{N})\neq\emptyset$, $A\cap(\{n_{2}\}\times\mathbb{N}^{<\omega})\neq\emptyset$.

It is similar to the case (2) and hence $A$ still has an upper bound in case (3).

(4) $A\cap(\{n_{0}\}\times\mathbb{N})$ is infinite and $A\cap(\{n_{1}\}\times\mathbb{N}^{<\omega})\neq\emptyset$, $A\cap(\{n_{2}\}\times\mathbb{N}^{<\omega})\neq\emptyset$.

Choose a $(n_{1},\mu_{1}\mu_{2}\cdot\cdot\cdot\mu_{k})\in(\{n_{1}\}\times\mathbb{N}^{<\omega})$ and $(n_{2},\nu_{1}\nu_{2}\cdot\cdot\cdot\nu_{k'})\in(\{n_{2}\}\times\mathbb{N}^{<\omega})$. Since $A\cap(\{n_{0}\}\times\mathbb{N})$ is infinite, there is a $(n_{0},\overline{m}_{0})\in A$ such that $$max\{\mu_{1}+2,\mu_{2}+2,\cdot\cdot\cdot,\mu_{k}+2,\nu_{1}+2,\nu_{2}+2,\cdot\cdot\cdot,\nu_{k'}+2\}<\overline{m}_{0}.$$
Now, $$\{(n_{0},\overline{m}_{0}),(n_{1},\mu_{1}\mu_{2}\cdot\cdot\cdot\mu_{k}),(n_{2},\nu_{1}\nu_{2}\cdot\cdot\cdot\nu_{k'})\}$$ has an upper bound $u$. Clearly, $u=(n_{1},\top)$ or $u=(n_{2},\top)$. we just consider the following two possibilities (4.1) and (4.2).

(4.1) $u=(n_{1},\top)$.

In this case, $(n_{2},\nu_{1}\nu_{2}\cdot\cdot\cdot\nu_{k'})\prec_{2}(n_{1},\top)$ and $(n_{0},\overline{m}_{0})\leq(n_{1},\top)$. By $(n_{2},\nu_{1}\nu_{2}\cdot\cdot\cdot\nu_{k'})\prec_{2}(n_{1},\top)$, $n_{1}=f_{n_{2}s_{1}}(\sigma_{1}\sigma_{2}\cdot\cdot\cdot\sigma_{k_{1}})$ for some $n_{2}<s_{1}$ and $\nu_{1}\nu_{2}\cdot\cdot\cdot\nu_{k'}\leq\sigma_{1}\sigma_{2}\cdot\cdot\cdot\sigma_{k_{1}}$.

If $k_{1}=1$, then $\nu_{1}\nu_{2}\cdot\cdot\cdot\nu_{k'}=\sigma_{1}\sigma_{2}\cdot\cdot\cdot\sigma_{k_{1}}=\nu_{1}$ and $(n_{0},\overline{m}_{0})\prec_{3}(f_{n_{2}s_{1}}(\nu_{1}),\top)$. Whence $\overline{m}_{0}<\nu_{1}$, impossible.

If $1<k_{1}$, then $(n_{0},\overline{m}_{0})\prec_{4}(f_{n_{2}s_{1}}(\sigma_{1}\sigma_{2}\cdot\cdot\cdot\sigma_{k_{1}}),\top)$. Hence, $n_{0}=f_{n_{2}s_{1}}(\sigma_{1}\sigma_{2}\cdot\cdot\cdot\sigma_{k_{1}-1})$ and $\overline{m}_{0}\leq\sigma_{k_{1}}$. Since $A\cap(\{n_{0}\}\times\mathbb{N})$ is infinite, there is a $(n_{0},\widetilde{m}_{0})\in A$ such that $$max\{\overline{m}_{0}+2,\sigma_{1}+2,\sigma_{2}+2,\cdot\cdot\cdot,\sigma_{k_{1}}+2\}<\widetilde{m}_{0}.$$
Then $\{(n_{0},\widetilde{m}_{0}),(n_{1},\mu_{1}\mu_{2}\cdot\cdot\cdot\mu_{k}),(n_{2},\nu_{1}\nu_{2}\cdot\cdot\cdot\nu_{k'})\}$ has an upper bound $u'$. Similarly, $u'=(n_{1},\top)$ or $u'=(n_{2},\top)$. Assume that $u'=(n_{2},\top)$. Then $(n_{1},\mu_{1}\mu_{2}\cdot\cdot\cdot\mu_{k})\prec_{2}u'=(n_{2},\top)$ and hence $n_{1}< n_{2}$ by strict order $\prec_{2}$. However, $n_{1}=f_{n_{2}s_{1}}(\sigma_{1}\sigma_{2}\cdot\cdot\cdot\sigma_{k_{1}})>n_{2}$, impossible. So we have $u'=(n_{1},\top)$. As $1<k_{1}$, $(n_{0},\widetilde{m}_{0})\prec_{4}(n_{1},\top)=(f_{n_{2}s_{1}}(\sigma_{1}\sigma_{2}\cdot\cdot\cdot\sigma_{k_{1}}),\top)$. This means that $\widetilde{m}_{0}\leq\sigma_{k_{1}}$, a contradiction.

(4.2) $u=(n_{2},\top)$.

It is similar to (4.1) and hence (4.2) is still impossible.

In a conclusion, the case (4) is always impossible.

(5) $A\cap(\{n_{0}\}\times\mathbb{N}^{<\omega})$ is infinite and $A\cap(\{n_{1}\}\times\mathbb{N}^{<\omega})\neq\emptyset$, $A\cap(\{n_{2}\}\times\mathbb{N}^{<\omega})\neq\emptyset$.

This case is always impossible since $n_{0},n_{1},n_{2}$ are different.

(6) $A\cap(\{n_{0}\}\times\mathbb{N}^{<\omega})$ is infinite and $A\cap(\{n_{1}\}\times\mathbb{N})\neq\emptyset$, $A\cap(\{n_{2}\}\times\mathbb{N}^{<\omega})\neq\emptyset$.

The details can see Appendix.


(7) $A\cap(\{n_{0}\}\times\mathbb{N}^{<\omega})$ is infinite and $A\cap(\{n_{1}\}\times\mathbb{N}^{<\omega})\neq\emptyset$, $A\cap(\{n_{2}\}\times\mathbb{N})\neq\emptyset$.

It is similar to case (6) and hence $A$ still has an upper bound in case (7).

(8) $A\cap(\{n_{0}\}\times\mathbb{N}^{<\omega})$ is infinite and $A\cap(\{n_{1}\}\times\mathbb{N})\neq\emptyset$, $A\cap(\{n_{2}\}\times\mathbb{N})\neq\emptyset$.

The details can see Appendix.

According to case (1)-(8), $A$ has an upper bound in $\widehat{P}$.\end{proof}

\end{document}
